% unidades/08-compras.tex
\chapter{🛍️ 買い物 — Compras}
\unidadheader{08}{買い物}{Compras}{Hablar de precios y elegir objetos}

\begin{tcolorbox}[vocabbox, title={📌 Vocabulario Nuevo}]
\begin{tabularx}{\linewidth}{llX}
\toprule
\textbf{Japonés} & \textbf{Español} & \textbf{Tipo} \\
\midrule
\vocabitem{いくら}{—}{cuánto (dinero)}{adv.}
\vocabitem{高い}{たかい}{caro / alto}{adj-い}
\vocabitem{安い}{やすい}{barato}{adj-い}
\vocabitem{大きい}{おおきい}{grande}{adj-い}
\vocabitem{小さい}{ちいさい}{pequeño}{adj-い}
\vocabitem{新しい}{あたらしい}{nuevo}{adj-い}
\vocabitem{古い}{ふるい}{viejo}{adj-い}
\vocabitem{鞄}{かばん}{bolsa / maleta}{sust.}
\vocabitem{靴}{くつ}{zapatos}{sust.}
\vocabitem{服}{ふく}{ropa}{sust.}
\bottomrule
\end{tabularx}
\end{tcolorbox}

\subsection*{🗣️ Frases Clave}

\frase{これはいくらですか？}{¿Cuánto cuesta esto?}
\frase{ちょっと\ruby{高}{たか}いですね。}{Es un poco caro, ¿verdad?}
\frase{もっと\ruby{安}{やす}いのはありますか？}{¿Hay algo más barato?}
\frase{これを\ruby{見}{み}せてください。}{Muéstreme esto, por favor.}

\subsection*{⚙️ Gramática}

\patron{Adjetivos con sustantivos}
  {Adj-い + Sustantivo}
  {\ej{\ruby{新}{あたら}しい\ruby{鞄}{かばん}を\ruby{買}{か}いました。}{Compré una bolsa nueva.}
  \ej{\ruby{小}{ちい}さい\ruby{靴}{くつ}が\ruby{好}{す}きです。}{Me gustan los zapatos pequeños.}}

\patron{Preguntar el precio}
  {〜はいくらですか？}
  {\ej{このカメラはいくらですか？}{¿Cuánto cuesta esta cámara?}
  \ej{\ruby{五万円}{ごまんえん}です。}{Son 50,000 yen.}}

\begin{tcolorbox}[ejerciciobox, title={📝 Ejercicios Rápidos}]
\begin{enumerate}
  \item ¿Cómo dices "Esta bolsa es barata"?
  \item Pregunta si hay zapatos más grandes.
  \item Traduce: この\ruby{服}{ふく}はとても\ruby{高}{たか}いです。
\end{enumerate}
\vspace{0.4em}
\textbf{Respuestas:}
1) このかばんは\ruby{安}{やす}いです。\quad
2) もっと\ruby{大}{おお}きい\ruby{靴}{くつ}はありますか？\quad
3) Esta ropa es muy cara.
\end{tcolorbox}

\begin{tcolorbox}[kanjibox, title={🀄 Kanjis de la Unidad}]
\begin{tabularx}{\linewidth}{cllXX}
\toprule
\textbf{Kanji} & \textbf{On} & \textbf{Kun} & \textbf{Significado} & \textbf{Vocab} \\
\midrule
\kanjirow{安}{あん}{やす}{barato / tranquilo}{\ruby{安い}{やすい} / \ruby{安心}{あんしん}}
\kanjirow{高}{こう}{たか}{caro / alto}{\ruby{高い}{たかい} / \ruby{高校}{こうこう}}
\kanjirow{大}{だい・たい}{おお}{grande}{\ruby{大きい}{おおきい} / \ruby{大学}{だいがく}}
\kanjirow{小}{しょう}{ちい}{pequeño}{\ruby{小さい}{ちいさい} / \ruby{小学校}{しょうがっこう}}
\kanjirow{色}{しょく・しき}{いろ}{color}{\ruby{色}{いろ} / \ruby{黄色}{きいろ}}
\bottomrule
\end{tabularx}
\end{tcolorbox}
